% GENERATED FILE. Edit canonical lessons, then run npm run generate:books.
% canonical-source-hash: fnv1a64:ea73b443c32cd784
% canonical-lessons: IT-C313-morte, IT-C313-defunto, IT-C313-funerale, IT-C313-lutto, IT-C313-seppellire

\chapter{L'ultimo saluto}
\label{ch:it-l-ultimo-saluto}
\hlchaptermodality{313}

\begin{chapteropening}
\textbf{By the end of this chapter:} \emph{I can talk about a death: the deceased, the funeral, mourning and burial.}

Complete the last lesson of chapter 313: I can talk about a death: the deceased, the funeral, mourning and burial.
\end{chapteropening}

\section[la morte]{la morte — death}

\label{lesson:IT-C313-morte}



\begin{quote}
Before the new one: say the Italian for the lineage, then the Italian for the family tree.
\end{quote}

\subsection*{You'll want to know: la morte}
\textbf{la morte} — \textquotedblleft{}death\textquotedblright{}.

\begin{quote}
La morte del nonno ci ha cambiato la vita.
\end{quote}

\begin{quote}
\emph{Grandfather's death changed our lives.}
\end{quote}

\emph{Annoiarsi a morte} — to be bored to death.

\begin{cousinweb}
\emph{morte} is inherited from Latin \emph{mors, mortis} \textquotedblleft{}death\textquotedblright{} (EtymDB).
\end{cousinweb}

\begin{grammarlens}[title={What you've built}]
One more word for telling, explaining and describing at B1.
\end{grammarlens}

\subsection*{Your turn}
\begin{itemize}
\raggedright
  \item \emph{Say it:} \emph{la morte}
  \item \emph{Say it:} \emph{la morte}, once more
  \item \emph{Say it:} a sentence of your own with \emph{la morte}
  \item \emph{Recall:} say the Italian for the lineage, then the Italian for the family tree, then say \emph{la morte} again
\end{itemize}

\begin{tcolorbox}[breakable,colback=teal!4,colframe=teal!35!black,title={Before you move on}]
What does la morte mean? (\textquotedblleft{}Death\textquotedblright{}.) Say it once more.
\end{tcolorbox}
\section[il defunto]{il defunto — the deceased}

\label{lesson:IT-C313-defunto}



\begin{quote}
Before the new one: say the Italian for the family tree, then the Italian for death.
\end{quote}

\subsection*{You'll want to know: il defunto}
\textbf{il defunto} — \textquotedblleft{}the deceased\textquotedblright{}.

\begin{quote}
La famiglia del defunto ringrazia tutti.
\end{quote}

\begin{quote}
\emph{The family of the deceased thanks everyone.}
\end{quote}

Also an adjective: \emph{il defunto marito}.

\begin{cousinweb}
\emph{defunto} is a learned borrowing from Latin \emph{dēfūnctus} \textquotedblleft{}one who has finished (life)\textquotedblright{} (EtymDB), from \emph{dēfungī} \textquotedblleft{}to discharge, to finish\textquotedblright{}.
\end{cousinweb}

\begin{grammarlens}[title={What you've built}]
One more word for telling, explaining and describing at B1.
\end{grammarlens}

\subsection*{Your turn}
\begin{itemize}
\raggedright
  \item \emph{Say it:} \emph{il defunto}
  \item \emph{Say it:} \emph{il defunto}, once more
  \item \emph{Say it:} a sentence of your own with \emph{il defunto}
  \item \emph{Recall:} say the Italian for the family tree, then the Italian for death, then say \emph{il defunto} again
\end{itemize}

\begin{tcolorbox}[breakable,colback=teal!4,colframe=teal!35!black,title={Before you move on}]
What does il defunto mean? (\textquotedblleft{}The deceased\textquotedblright{}.) Say it once more.
\end{tcolorbox}
\section[il funerale]{il funerale — the funeral}

\label{lesson:IT-C313-funerale}



\begin{quote}
Before the new one: say the Italian for death, then the Italian for the deceased.
\end{quote}

\subsection*{You'll want to know: il funerale}
\textbf{il funerale} — \textquotedblleft{}the funeral\textquotedblright{}.

\begin{quote}
Al funerale c'era tutto il paese.
\end{quote}

\begin{quote}
\emph{The whole village was at the funeral.}
\end{quote}

Often plural: \emph{i funerali}.

\begin{cousinweb}
\emph{funerale} is inherited from Late Latin \emph{fūnerālis} (EtymDB), built on \emph{fūnus} \textquotedblleft{}a burial\textquotedblright{}.
\end{cousinweb}

\begin{grammarlens}[title={What you've built}]
One more word for telling, explaining and describing at B1.
\end{grammarlens}

\subsection*{Your turn}
\begin{itemize}
\raggedright
  \item \emph{Say it:} \emph{il funerale}
  \item \emph{Say it:} \emph{il funerale}, once more
  \item \emph{Say it:} a sentence of your own with \emph{il funerale}
  \item \emph{Recall:} say the Italian for death, then the Italian for the deceased, then say \emph{il funerale} again
\end{itemize}

\begin{tcolorbox}[breakable,colback=teal!4,colframe=teal!35!black,title={Before you move on}]
What does il funerale mean? (\textquotedblleft{}The funeral\textquotedblright{}.) Say it once more.
\end{tcolorbox}
\section[il lutto]{il lutto — mourning, bereavement}

\label{lesson:IT-C313-lutto}



\begin{quote}
Before the new one: say the Italian for the deceased, then the Italian for the funeral.
\end{quote}

\subsection*{You'll want to know: il lutto}
\textbf{il lutto} — \textquotedblleft{}mourning, bereavement\textquotedblright{}.

\begin{quote}
La città è in lutto.
\end{quote}

\begin{quote}
\emph{The city is in mourning.}
\end{quote}

\emph{Essere in lutto} — to be in mourning.

\begin{cousinweb}
\emph{lutto} is inherited from Latin \emph{lūctus} \textquotedblleft{}grief, sorrow\textquotedblright{} (EtymDB).
\end{cousinweb}

\begin{grammarlens}[title={What you've built}]
One more word for telling, explaining and describing at B1.
\end{grammarlens}

\subsection*{Your turn}
\begin{itemize}
\raggedright
  \item \emph{Say it:} \emph{il lutto}
  \item \emph{Say it:} \emph{il lutto}, once more
  \item \emph{Say it:} a sentence of your own with \emph{il lutto}
  \item \emph{Recall:} say the Italian for the deceased, then the Italian for the funeral, then say \emph{il lutto} again
\end{itemize}

\begin{tcolorbox}[breakable,colback=teal!4,colframe=teal!35!black,title={Before you move on}]
What does il lutto mean? (\textquotedblleft{}Mourning, bereavement\textquotedblright{}.) Say it once more.
\end{tcolorbox}
\section[seppellire]{seppellire — to bury}

\label{lesson:IT-C313-seppellire}



\begin{quote}
Before the new one: say the Italian for the funeral, then the Italian for mourning.
\end{quote}

\subsection*{You'll want to know: seppellire}
\textbf{seppellire} — \textquotedblleft{}to bury\textquotedblright{}.

\begin{quote}
Lo hanno sepolto nel paese dove era nato.
\end{quote}

\begin{quote}
\emph{They buried him in the village where he was born.}
\end{quote}

Past participle \emph{sepolto}, also \emph{seppellito}.

\begin{cousinweb}
\emph{seppellire} is inherited from Latin \emph{sepelīre} \textquotedblleft{}to bury\textquotedblright{} (Lewis \& Short).
\end{cousinweb}

\begin{grammarlens}[title={What you've built}]
One more word for telling, explaining and describing at B1.
\end{grammarlens}

\subsection*{Your turn}
\begin{itemize}
\raggedright
  \item \emph{Say it:} \emph{seppellire}
  \item \emph{Say it:} \emph{seppellire}, once more
  \item \emph{Say it:} a sentence of your own with \emph{seppellire}
  \item \emph{Recall:} say the Italian for the funeral, then the Italian for mourning, then say \emph{seppellire} again
\end{itemize}

\begin{tcolorbox}[breakable,colback=teal!4,colframe=teal!35!black,title={Before you move on}]
What does seppellire mean? (\textquotedblleft{}To bury\textquotedblright{}.) Say it once more.
\end{tcolorbox}
